%% ============================================================
%%  Institutional Background and Identification Strategy
%%  Depositor Discipline in Modern Banking (2026-04-01 draft)
%% ============================================================

\section{Institutional Background and Identification Strategy}
\label{sec:background}

\subsection{Credit-Loss Accounting Before CECL}
\label{sec:background:ilm}

For most of the post-war period, U.S. banks recognized credit losses under the
\textit{incurred loss model} (ILM), codified in Accounting Standards Codification
(ASC) 450-20 and ASC 310-10. Under the ILM, a bank was prohibited from booking a
loss reserve until a triggering event—typically a missed payment or an internal
classification downgrade—made future default ``probable and estimable.'' Because
the standard anchored provisioning to observable credit events that had already
occurred, it systematically deferred recognition of the latent risk embedded in
seasoning loan portfolios.\footnote{The procyclical bias of the incurred-loss
standard has long been noted in the literature. Provisioning lags default rates
by one to two years over the cycle \citep[e.g.,][]{BeattyLiao2011}, allowing
banks to report inflated capital ratios through downturns before the full cost of
earlier lending vintages is acknowledged.} As a result, two otherwise identical
banks could carry materially different loan-loss reserves depending solely on
idiosyncratic differences in portfolio vintage, borrower geography, and the
timing of internal credit reviews—none of which necessarily reflected differences
in \textit{current} managerial risk appetite.

The 2007–2009 financial crisis crystallized regulatory concern that the ILM was
masking systemic fragility. In its aftermath, the Financial Accounting Standards
Board (FASB) launched a multi-year project to replace the incurred-loss trigger
with a forward-looking recognition model.

\subsection{The CECL Standard}
\label{sec:background:cecl}

In June 2016, FASB issued Accounting Standards Update (ASU) 2016-13, which
replaced the incurred-loss framework with the \textit{Current Expected Credit
Loss} standard (CECL), codified in ASC 326. Under CECL, a bank must estimate and
reserve for \textit{lifetime} expected credit losses at the moment a financial
asset is originated or acquired. The reserve estimate must incorporate historical
loss experience, current macroeconomic conditions, and forward-looking forecasts.
No specific default trigger is required: a bank that originates a one-hundred-dollar
loan on day one must immediately provision an expected loss spanning the full
contractual life of that loan.

The economic magnitude of the required reserve is a direct function of the
risk composition of the loan portfolio at the adoption date. Loans with longer
remaining maturities, weaker borrower credit quality, or underrepresented collateral
positions require larger CECL reserves relative to the ILM allowance they
previously carried. Consequently, the difference between the new CECL-required
allowance and the existing allowance for loan and lease losses (ALLL) accumulated
under the ILM varies substantially across banks, driven by the cumulative
lending decisions of years prior to adoption.

\paragraph{Adoption timeline.}
FASB phased adoption in by entity type. Large public companies (SEC registrants
with fiscal years beginning after December 15, 2019) were required to adopt CECL
on January 1, 2020. Smaller reporting companies, non-accelerated filers, and
non-public business entities—which constitute our main sample of community
banks—were required to adopt on January 1, 2023. Our identification therefore
centers on the community-bank cohort, with the earlier large-bank adoption
providing a signal-validation sample free from the regulatory confounds
introduced by the COVID-19 shock (Section~\ref{sec:signal}).

\paragraph{Regulatory capital phase-in.}
Concurrent with the 2023 adoption mandate, banking regulators permitted community
banks to elect a three-year phase-in of the day-one regulatory capital impact.
Under this election, a bank phases in the reduction to risk-weighted capital
ratios over 2023–2025 rather than absorbing it immediately. Critically, the
phase-in affected only the \textit{regulatory capital} treatment of the shock, not
its accounting recognition on the public financial statements. A bank electing the
phase-in still reported the full CECL reserve and the full retained-earnings
charge in its Q1 2023 Call Report. We exploit variation in phase-in election as
a heterogeneity test in Section~\ref{sec:heterogeneity}.

\subsection{The Day-One Adjustment}
\label{sec:background:dayone}

Upon adoption, each bank computed the difference between the CECL-required
allowance and its existing ILM allowance. This difference was recorded as a
cumulative-effect adjustment directly to retained earnings on the opening balance
sheet of the adoption quarter—what practitioners term the ``Day-One'' impact.
Formally, for bank $i$ adopting in quarter $t^{*}$:
\begin{equation}
  \Delta_i \;\equiv\; \underbrace{\mathrm{ALLL}^{\mathrm{CECL}}_{i,t^{*}}}_{\text{new allowance}}
                    - \underbrace{\mathrm{ALLL}^{\mathrm{ILM}}_{i,t^{*}-1}}_{\text{existing allowance}},
  \label{eq:dayone}
\end{equation}
where $\Delta_i < 0$ for the vast majority of adopters (the new allowance exceeds
the existing one, implying a retained-earnings charge). Our treatment variable
scales this adjustment by beginning-of-period equity to ensure comparability
across banks of different sizes:
\begin{equation}
  \texttt{cecl\_equity}_{i}
  \;=\;
  \frac{\Delta_i}{\texttt{equity\_bop}_{i,t^{*}}}
  \times 100,
  \label{eq:treatment}
\end{equation}
expressed in percentage points. A value of $-5$, for example, indicates that
the CECL charge erased five percentage points of the bank's pre-adoption equity.
Because we want treatment intensity to reflect the ex ante CECL shock rather than
post-adoption accounting movements, we assign each bank its adoption-quarter
value of \texttt{cecl\_equity} and hold it fixed throughout the event-time
panel.\footnote{In a small number of cases, banks report their CECL adjustment in
a quarter adjacent to the nominal 2023Q1 date (due to fiscal-year heterogeneity
or filing lags). Our adoption-date algorithm identifies the first quarter in the
2022Q3–2023Q4 rollout window with a non-missing \texttt{cecl\_equity} observation,
and assigns that quarter as the bank-specific event time $t^{*}=0$.}

For our main binary specification we classify a bank as \textit{high-CECL} if
its adoption-quarter shock exceeds a pre-specified threshold:
\begin{equation}
  \texttt{high\_cecl\_equity}_{i}
  \;=\;
  \mathbf{1}\!\left\{
    \texttt{cecl\_equity}_{i} \;\geq\; \tau
  \right\},
  \label{eq:binary}
\end{equation}
where $\tau$ corresponds to the top-decile cutoff of the cross-sectional CECL
distribution in the community-bank sample.\footnote{The threshold is set at
\texttt{CECL\_THRESHOLD\_PCT} in our analysis code. All results are robust to
alternative threshold choices, including continuous-treatment specifications
that replace the binary indicator with the raw \texttt{cecl\_equity} value.}

\subsection{Identification Strategy}
\label{sec:background:id}

Our goal is to identify the causal effect of disclosed credit-risk information on
depositor behavior. The standard challenge is that risk and funding structure are
jointly determined: riskier banks attract fewer uninsured deposits at any point in
time, so cross-sectional comparisons conflate information responses with
equilibrium sorting. We argue that the CECL Day-One adjustment provides a rare
setting in which risk-relevant information arrives at a \textit{known moment} and
with \textit{ex ante predetermined magnitude}, allowing us to separate the
informational content of a disclosure from concurrent endogenous changes in bank
policy.

We ground our design in three conditions.

\paragraph{(i) Timing is exogenous.}
The adoption date for community banks was fixed by FASB and federal banking
regulators, not chosen by individual banks. A bank could not delay adoption to
avoid revealing an unfavorable Day-One charge, nor could it accelerate adoption to
pre-empt depositor scrutiny. Absent cross-bank lobbying that differentially shifted
the mandate date, the \textit{when} of the shock is orthogonal to bank-level
characteristics as of 2022.

\paragraph{(ii) Magnitude is predetermined by historical lending decisions.}
The size of $\Delta_i$ in \eqref{eq:dayone} reflects the cumulative composition
of a bank's loan portfolio as it stood in years prior to adoption—loan vintages,
maturities, collateral coverage, and borrower quality accumulated well before the
standard was finalized. While management can respond to anticipated CECL exposure
by gradually tightening underwriting after 2016, such responses compress, rather
than eliminate, cross-sectional variation in $\Delta_i$; they do not generate
endogenous sorting that would bias our estimates. Moreover, the portion of
variation in \texttt{cecl\_equity} attributable to portfolio vintage is unlikely
to correlate with contemporaneous depositor behavior through any channel other than
the information disclosed in the quarterly Call Report filing.

\paragraph{(iii) The shock is informational, not liquidity-driven.}
Equation~\eqref{eq:dayone} records an accounting charge with no immediate cash
consequence. The Day-One adjustment does not trigger a required cash transfer,
dividend cut, or borrower default. It reduces regulatory capital ratios (unless the
bank elects phase-in) and reported book equity, but it does not alter the cash
flows, contractual terms, or physical performance of the underlying loans.
Accordingly, any change in deposit quantities or deposit pricing that we observe
after the adoption quarter cannot be attributed to a mechanical liquidity shock or
a contemporaneous deterioration in underlying credit performance. It must instead
reflect depositor reaction to the \textit{information} released by the Day-One
disclosure.

This separation of information from cash flows is the central advantage of our
setting relative to alternative tests of depositor discipline based on realized
loan defaults, dividend omissions, or rating downgrades, each of which confounds
new information with contemporaneous changes in bank liquidity or cash position.

\paragraph{Empirical design.}
We implement a difference-in-differences design that exploits both the
\textit{time-series} variation (before versus after the 2023Q1 adoption quarter)
and the \textit{cross-sectional} variation in shock magnitude. Our baseline
specification is:
\begin{equation}
  y_{it}
  \;=\;
  \alpha_i + \lambda_t
  + \beta_1 \cdot \texttt{post}_{it}
    \times \texttt{high\_cecl\_equity}_{i}
  + X_{i,t-1}'\gamma
  + \varepsilon_{it},
  \label{eq:baseline}
\end{equation}
where $y_{it}$ is a deposit-composition or funding-cost outcome for bank $i$ in
quarter $t$; $\alpha_i$ are bank fixed effects; $\lambda_t$ are calendar-quarter
fixed effects; $\texttt{post}_{it}$ equals one in quarters at or after the
bank's adoption date; and $X_{i,t-1}$ is a vector of lagged controls including
log total assets and equity-to-assets.\footnote{We also estimate continuous-treatment
variants that replace the binary interaction with $\texttt{post}_{it}\times
\texttt{cecl\_equity}_i$, and dynamic event-study specifications that replace the
post indicator with a full set of quarter-relative-to-adoption dummies over the
window $|\texttt{qtrs\_since}| \leq 10$, with $k=-1$ as the reference period.
Identification in the dynamic specification requires that high- and low-CECL banks
trend in parallel prior to adoption—a condition we evaluate empirically in
Section~\ref{sec:results}.}

The identifying assumption is that, conditional on bank and time fixed effects and
lagged controls, the CECL shock magnitude is uncorrelated with differential trends
in depositor behavior that would have occurred absent adoption. Parallel pre-trends
in event-study plots provide our primary diagnostic. We complement this check with
a signal-validation exercise using large banks—which adopted CECL three years
earlier under otherwise similar conditions—to confirm that equity markets and credit
derivatives priced the CECL information shock at the time of adoption, supporting
its interpretation as a genuine disclosure event (Section~\ref{sec:signal}).